\problemname{Hidden Permutation}

When a fixed shuffle is applied to a deck of cards repeatedly, the deck will eventually return to its original order. 
The number of times the shuffle must be applied before this happens is called the \textit{period} of the deck for that shuffle.
In the present task, we consider the inverse question: rather than starting with a shuffle and finding its period, 
you are told, for each possible period, how many decks of cards have that period. Your task is to construct any 
shuffle consistent with this information.

Formally, there is a hidden permutation $f = (f_1, f_2, \dots, f_N)$, where $N$ is a given positive integer. A permutation is a list of numbers where
each number from $1$ to $N$ occurs exactly once. Let $x = x_1x_2 \dots x_N$ be a binary string. We say that $f(x)$ ($f$ applied to $x$) is the new
binary string $x_{f_1}x_{f_2} \dots x_{f_N}$, and the \textit{period} of $x$ is the smallest positive integer $p$ such that 

$$\underbrace{f(f(\cdots f(x)\cdots))}_{p\ \text{times}} =  x.$$

It can be proven that for any permutation $f$ and any binary string $x$ of the same length, the period of $x$ exists. In other words, 
if you apply $f$ to $x$ enough times, you eventually get back $x$.

You are given the number $N$, and for each possible period $p$ you are given the number of binary strings 
 (out of all $2^N$ possible binary strings) that have period $p$. Your task is to find any permutation $f$ that corresponds to the given information, 
 or conclude that there is no such permutation.

\begin{Input}
    The input consists of:
    \begin{itemize}
      \item One line with integers $N$, $K$ ($2 \leq N \leq 100$, $1 \leq K \leq 1000$), the length of the permutation and the number of possible periods.

      \item One line with $K$ integers $p_1, p_2, \dots, p_K$ ($1 \leq p_i \leq 10^9$). These are all the possible periods that binary strings of 
      length $N$ can have, with respect to the hidden permutation. The numbers $p_i$ are all distinct, and are sorted in increasing order.

      \item One line with $K$ integers $m_1, m_2, \dots, m_K$ ($1 \leq m_i \leq 2^{100}$). This means 
      that there are $m_i$ binary strings with period $p_i$.

    \end{itemize}
    
    Note that the numbers $m_i$ can be quite large!

\end{Input}

\begin{Output}
    If there is no permutation $f$ that corresponds to the given information, print \verb|impossible|.
    Otherwise, print one line with $N$ integers $f_1, f_2,  \dots, f_N$, the permutation that you chose. Recall that the permutation 
    should contain every integer from $1$ to $N$ exactly once. If there are multiple solutions, you can print any one of them.
\end{Output}

\section*{Sample Explanation}

In the second sample, the input almost corresponds to the permutation $(2,3,4,\dots, 91, 1)$. The only error is that the last digit in the last number 
of the input should be $0$ instead of $1$.

In the third sample, all $4$ binary strings have period $1$, which means that the identity permutation $(1,2)$ is a valid answer.
